\subsection{A depth-three limitation at quintic primes}
\label{subsec:ff3-quintic-depth-three}

The full pro-three boundary theorem cannot in general be detected
in the first three ordinary lower-central grades.  The following
result gives a specified positive-density family and an explicit prime.

\begin{proposition}[Constant quintic point map through class three]
\label{prop:ff3-quintic-depth-three}
Let $q\ne3,5$ be a rational prime with $q\equiv1\pmod5$, let
\[
 C/\mathbb F_q:\ X^5+Y^5+Z^5=0,\qquad
 C_A:\ v^2+v=u^5,\qquad A=\operatorname{Jac}(C_A),
\]
and suppose arithmetic Frobenius on $A[3]$ has order twenty.
Use $b=(0:-1:1)$ as base point and put
\[
 N=\pi_1(C_{\overline{\mathbb F}_q},b)^{\mathrm{pro}\text{-}3}.
\]
With closed ordinary lower-central terms, continuous nonabelian
cohomology satisfies
\[
 H^1(G_{\mathbb F_q},N/\gamma_4N)=\{1\}.
\]
Thus every section of this arithmetic quotient is geometrically
conjugate to the point section of $b$.  In particular, the proper
point-section map is constant in this quotient and in every finite
Galois-stable quotient of $N$ of nilpotency class at most three.

The set of rational primes satisfying these conditions has natural
density $1/10$, and it contains $q=11$.
\end{proposition}

\begin{proof}
The exact residual calculation for the quintic CM Jacobian gives
\[
 E=\mathbb Q(\zeta_5)\subset K=\mathbb Q(A[3]),\qquad
 [K:\mathbb Q]=80,\qquad
 \operatorname{Gal}(K/E)=C_{20}\subset\mathbb F_{81}^{\times};
\]
see Lemma~\ref{lem:lm-quintic-residual-group}.
Here $A[3]$ has rank one over
$\mathbb Z[\zeta_5]/3=\mathbb F_{81}$.
The three degree-five quotient maps
\[
 \begin{split}
 f_Z&=(-XY/Z^2,X^5/Z^5),\\
 f_X&=(-YZ/X^2,Y^5/X^5),\\
 f_Y&=(-ZX/Y^2,Z^5/Y^5)
 \end{split}
\]
give the integral isomorphism
\[
 T_3\operatorname{Jac}(X_5)\simeq(T_3A)^3.
\]
Indeed their joint pushforward $\Phi$ and summed pullback $\Psi$
satisfy $\Phi\Psi=[5]$ and $\Psi\Phi=[5]$, by the three disjoint
diagonal-character summands.  Hence $5^{-1}\Psi$ is an integral
three-adic inverse.  Good reduction at $q$ transports this identity
to the finite-field Frobenius representation.

Because $q$ splits in $E$, its residual Frobenius is multiplication
by an element $t\in\mathbb F_{81}^{\times}$ of order twenty.
The eigenvalues over $\overline{\mathbb F}_3$ on $A[3]$ are
\[
 t,\quad t^3,\quad t^9,\quad t^{27}=t^7.
\]
Each occurs three times on
$\overline V=N^{\mathrm{ab}}/3$.
The action is semisimple since twenty is prime to three.
None of the exponents $1,3,7,9$ is zero modulo twenty;
every sum of two lies between two and eighteen;
every sum of three is odd.
Consequently $1$ is not an eigenvalue of Frobenius on
\[
 \overline V^{\otimes m},\qquad 1\le m\le3.
\]
Repeated tensor factors are included.  Replacing $t$ by another
generator, or Frobenius by its inverse, does not change this assertion.

Put $L_m=\gamma_mN/\gamma_{m+1}N$.
Iterated commutators give a continuous equivariant surjection
\[
 (N^{\mathrm{ab}})^{\otimes m}\longrightarrow L_m;
\]
the grades are finitely generated compact $\mathbb Z_3$-modules.
The commutator description and its compatibility with the Galois
action are also described in
\citet[Introduction, author pp.~1--2]{DavisPriesWickelgren2024}.
It follows that $F-1$ is invertible modulo three on $L_m$ for
$m\le3$. Nakayama's lemma makes it surjective integrally, and a
surjective endomorphism of a finitely generated $\mathbb Z_3$-module
is injective. Thus
\[
 H^0(G_{\mathbb F_q},L_m)=H^1(G_{\mathbb F_q},L_m)=0
 \qquad(1\le m\le3).
\]
For $H^1$, explicitly choose a coboundary with the prescribed
cocycle value at the topological Frobenius generator; equality on
the full procyclic group follows by continuity.  The arithmetic
quotient remains full $G_{\mathbb F_q}$.

Induct through the central extensions
\[
 1\longrightarrow L_m\longrightarrow N/\gamma_{m+1}N
 \longrightarrow N/\gamma_mN\longrightarrow1.
\]
The abelian case follows from the displayed vanishing.  In a
successive stage, conjugate a cocycle so that its image in the
preceding quotient is trivial, lifting one conjugating element
through the surjection.  The remaining cocycle takes values in
$L_m$ and is a coboundary.  This proves the asserted singleton
cohomology set through $m=3$.  Equivalently, an invariant path lifts
at each central step because its fibre is an $L_m$-torsor with
zero $H^1$; the zero $H^0$ also makes these low-depth paths unique.
No section in the full geometric group is assumed to exist above
an arbitrary quotient section.

The residual elements defining the prime set are the eight generators
of the normal subgroup $C_{20}$ in the group of order eighty.
This is a conjugacy-stable subset, so Chebotarev gives natural
density $8/80=1/10$.

For an explicit prime, exact counting on $v^2+v=u^5$ gives
\[
 \#C_A(\mathbb F_{11})=13,\qquad
 \#C_A(\mathbb F_{121})=103.
\]
One can use $\mathbb F_{121}=\mathbb F_{11}[s]/(s^2-2)$,
where $2$ is a nonsquare, and count the square roots of
$1+4u^5$.  The Frobenius polynomial is therefore
\[
 T^4+T^3-9T^2+11T+121.
\]
Modulo three it is $T^4+T^3+2T+1$.
Its derivative is $(T-1)^3$ and its value at one is two, so it is
square-free.  Its companion matrix has order twenty by exact
matrix multiplication over $\mathbb F_3$.  Thus $q=11$ meets
the hypothesis.  These finite calculations are recorded in the
reproducibility certificate.
\end{proof}

At $q=11$, the same isogeny gives $\#C(\mathbb F_{11})=15$,
so all proper rational points are the fifteen coordinate points.
Their equality in $N/\gamma_4N$ is compatible with
Theorem~\ref{thm:ff3-different-components}: the latter concerns
the full geometric pro-three group and separates different
coordinate components. The proposition makes no assertion of a
full collision within one component. Already in tensor degree four,
$1+1+9+9\equiv0\pmod{20}$, so the eigenvalue exclusion stops.
An actual integral higher-path class remains to be determined;
this limitation of low-depth quotients does not settle the affine
base-five or base-seven Selmer-section problem.
